\section{Controlling what the ranker can see}
\label{sec:visibility}

\subsection{Setting and notation}
\label{sec:setting}

A scan is a set of native voxel centres $v\in\mathbb{R}^3$ (in millimetres, physical coordinates) with intensities $I(v)$ in Hounsfield units (HU). The candidate generator supplies proposals $p\in\mathbb{R}^3$; we use the released LUNA16 proposals unchanged~\cite{setio2017luna16}. For each proposal, a ranker receives a crop and returns a score. The crop is a $32^3$ grid centred at $p$ and aligned with the native axes, with $1.2\mm$ spacing, obtained by trilinear interpolation of the native voxels. Values are clipped to $[-1000,400]\HU$ and stored as 8-bit integers; the networks use the central $28^3$ samples, which span $32.4\mm$. All conditions below share this crop, so the input shape never changes.

\subsection{Mask native voxels, then resample}
\label{sec:mask}

A \emph{visibility condition} is a set $V$ of native voxels the ranker may see. All other voxels are replaced by air:
\begin{equation}
 I_V(v)=
 \begin{cases}
 I(v), & v\in V,\\
 -1000\HU, & v\notin V .
 \end{cases}
 \label{eq:mask}
\end{equation}
The crop is then resampled from $I_V$. The order matters. Trilinear interpolation mixes up to eight neighbouring native voxels, so masking the \emph{resampled} grid would let a hidden voxel near the boundary leak into a visible sample. With Eq.~\eqref{eq:mask}, a hidden voxel contributes only the constant value of air (Fig.~\ref{fig:visibility}b). \paragraph{Implementation checks.} For speed, masking is applied to a cropped neighbourhood of each proposal rather than to the whole scan, with a margin of native voxels kept for interpolation. We compared this implementation with literal whole-volume masking in nine test settings that include rotated, reflected and anisotropic native grids and proposals near the scan boundary. The crops were bit-identical. Replacing all 485,801 hidden voxel values in these settings with random values also left every crop and every sparse sample unchanged, and all 15,664 positive-weight interpolation contributions fell inside the visible set. The network receives no separate visibility channel; hidden regions appear only as air (supplement Sec.~S2).

The three main conditions use the ball
\begin{equation}
B(p)=\{v:\lVert v-p\rVert\le 10\mm\},
\end{equation}
the ball plus a band $A(p)$ defined next, and the full field.

\subsection{The band: where sparse samples get their values}
\label{sec:band}

Let $d$ be the diagonal of one native voxel, which depends on the scan's spacing. Any native voxel with a positive trilinear weight for a sample point $q$ lies within $d$ of $q$. We place sparse samples at two radii from the proposal,
\begin{equation}
 \begin{aligned}
 r_N &=10\mm+\max(1.9\mm,\,d+0.05\mm),\\
 r_F &=r_N+1.5\mm .
 \end{aligned}
 \label{eq:radii}
\end{equation}
By the triangle inequality, every native voxel $v$ that contributes to a sample at either radius satisfies
\begin{equation}
 r_N-d\;\leq\;\lVert v-p\rVert\;\leq\;r_F+d .
 \label{eq:support}
\end{equation}
The band is exactly this shell of native voxels,
\begin{equation}
A(p)=\{v:\ r_N-d\le\lVert v-p\rVert\le r_F+d\}.
\label{eq:band}
\end{equation}
Because $r_N-d\ge10.05\mm$, the band never overlaps the ball, and a thin gap of air separates them. For example, with $0.70\times0.70\times1.25\mm$ voxels, $d=1.59\mm$, $r_N=11.9\mm$, $r_F=13.4\mm$, and the band covers $10.3$--$15.0\mm$. The construction adapts to each scan's spacing, so Eq.~\eqref{eq:support} holds for every slice thickness. Algorithm~\ref{alg:band} summarizes the procedure.

\begin{algorithm}[t]
\caption{Building one visibility-controlled input}
\label{alg:band}
\small
\begin{algorithmic}[1]
\Require scan $\{(v,I(v))\}$, voxel spacing $s$, proposal $p$, condition $c\in\{\text{ball},\text{band},\text{sources},\text{full}\}$
\State $d\gets\lVert s\rVert_2$ \Comment{native voxel diagonal}
\State $r_N\gets 10+\max(1.9,\,d+0.05)$;\ \ $r_F\gets r_N+1.5$ \Comment{Eq.~\eqref{eq:radii}}
\State $B\gets\{v:\lVert v-p\rVert\le10\}$
\State $A\gets\{v:r_N-d\le\lVert v-p\rVert\le r_F+d\}$ \Comment{Eq.~\eqref{eq:band}}
\State $U\gets\{v:\ v$ has positive trilinear weight for one of the 384 samples $q_{j,\pm,r}\}$
\State $V\gets B$, $B\cup A$, $B\cup U$ or all voxels, according to $c$
\State $I_V(v)\gets I(v)$ if $v\in V$, else $-1000$ \Comment{Eq.~\eqref{eq:mask}}
\State $X\gets$ trilinear resample of $I_V$ on the $32^3$ grid at $1.2\mm$ around $p$
\State \Return $\mathrm{uint8}(\mathrm{clip}(X,-1000,400))$
\end{algorithmic}
\end{algorithm}

\subsection{Sparse samples and their native sources}
\label{sec:source_method}

We choose 96 unit axes $u_1,\dots,u_{96}$ spread evenly over a hemisphere (a Fibonacci lattice). On each axis we sample both directions ($\pm$) at both radii, giving $96\times2\times2=384$ samples $q_{j,\pm,r}=p\pm r\,u_j$ with $r\in\{r_N,r_F\}$. All sparse readouts in Sec.~\ref{sec:sparse_method} receive the same 384 clipped and normalized values, plus a background \emph{floor}: for each axis, the median of the samples on all other axes. The band $A(p)$ is the envelope of the support of these samples.

The band differs from the sparse samples in two ways at once: it shows far more voxels, and it enters the image network before the first convolution, whereas the sparse samples enter a separate head. To separate these, we add a fourth image condition. Let $U(p)$ be the set of native voxels with a strictly positive trilinear weight for at least one of the 384 samples. The \emph{native-source} condition uses $V=B(p)\cup U(p)$ in Eq.~\eqref{eq:mask}, with the same two-scale network and training schedule. It gives the image network every native value that the sparse samples are built from, at the same point of entry as the band. On average, $U(p)$ contributes \QuerySources{} exterior voxels to the crop, against \BandSources{} for the band (Fig.~\ref{fig:coverage}a).

\subsection{Readouts of the sparse samples}
\label{sec:sparse_method}

The sparse heads test whether the 384 values are enough when read well. Each head starts from a ball-only image ranker that is already trained and frozen. Let $s$ be its score (logit). A head predicts a correction $h$ from the samples and the clipped score, and the final score is $s+h$. The last layer of $h$ is initialized to zero, so every head starts exactly at the ball-only ranker, which is its baseline.

\paragraph{Tokens.} A shared encoder maps the two radii on one arm, taken relative to the floor, to a vector. For axis $j$ this gives $a_j$ (positive arm) and $b_j$ (negative arm). The axis token is $[a_j+b_j,\ a_j\odot b_j,\ |a_j-b_j|]$, which does not depend on which arm is called positive.

\paragraph{Direction-blind and direction-aware attention.} Two attention blocks process the 96 tokens, followed by mean and max pooling. Let $z_{jk}=|u_j^\top u_k|\in[0,1]$ be the absolute cosine between axes $j$ and $k$. The \emph{direction-aware} head adds a fixed bias to the attention logits of its four heads: $12(z_{jk}-1)$ and $35(z_{jk}-1)$, which favour neighbouring axes at two scales, $-6z_{jk}$, which favours well-separated axes, and $0$ (global). The \emph{direction-blind} head uses zero bias in all four heads. Otherwise the two heads are identical: same initial weights, training batches, objective, checkpoints and parameter count. Tokens and the bias matrix are permuted together during training, so neither head can read direction from the token order. We train both at widths 48 and 96 (54k and 213k parameters).

\paragraph{Compact heads and controls.} A smaller head (873 parameters) pools the axis tokens without attention. We refit it with the same frozen scores and training data as the attention heads. Further controls test what such a compact head uses: summarizing each axis by one scalar instead of its endpoints, fusing an image embedding instead of the score, randomly re-pairing opposite arms (\emph{sham}), and removing the pairing. A last test adds the paired head on top of the full-field ranker. Details and all control results are in the supplement (Secs.~S2--S3).

\section{Experimental setup}
\label{sec:evaluation}

\paragraph{Data and roles.} We use the LUNA16 release of LIDC-IDRI, with its V2 candidate coordinates, reference nodules and excluded findings~\cite{setio2017luna16,armato2011lidc}. Scans are split by patient: 356 for training, 44 for checkpoint selection, 45 for development choices and 88 for threshold calibration. The evaluation pool contains \EvalScans{} scans from subsets 6--9, with \EvalMarks{} proposals and \EvalLesions{} reference lesions; at most \ProposedLesions{} lesions have a matching proposal. Two scans of one repeated subject are removed. Earlier related studies, including crop studies, used every nodule-bearing scan of this pool, so all results are retrospective evaluations on a reused benchmark. Each comparison and the anatomy grouping were fixed before evaluation scores were computed.

\paragraph{Rankers.} The main ranker is a \emph{two-scale} 3D CNN that reads the central high-resolution grid and a downsampled context grid. A \emph{single-stream} 3D network first rearranges $2^3$ blocks into channels and processes the full field in one residual hierarchy. A \emph{three-view} 2D model reads three orthogonal stacks of central slices. All rankers see the same masked crops. Every image condition is trained for 48 epochs with the same optimizer, augmentation and hard-negative mining. The checkpoint with the best CPM on the selection scans is kept. The two two-scale runs (\emph{fit~1} and \emph{fit~2}) differ only in their training seed.

\paragraph{Implementation details.} The two-scale network reads a central $16^3$ view and a $14^3$ average-pooled view of the $28^3$ input. Each stream has a $3^3$ stem with 16 channels and three stride-two stages of widths 32, 64 and 96; pooled mean and max features feed a 96-unit projection (1.28M parameters in total). The single-stream network turns each $2^3$ block into eight channels on a $14^3$ grid, followed by residual stages of widths 32, 64 and 96 (0.68M parameters). The three-view model reads six central slices per orientation with a shared 2D encoder. Each epoch has 160 batches of 24 positives and 104 negatives. Augmentation uses shifts of up to four grid steps, right-angle in-plane rotations, flips, and intensity gain and offset; the visibility mask is applied before augmentation. We use AdamW (learning rate 0.0018, cosine decay to 12\%, weight decay 0.012). Every fourth epoch, the 16,000 highest-scoring training negatives are refreshed, and half of each batch's negatives are drawn from this pool. Sparse heads train for 24 epochs on the same batches, with checkpoints every four epochs; ties are broken by the earliest epoch.

\paragraph{Metrics.} A proposal detects a lesion if it lies within the reference radius; only the highest-scoring detection per lesion counts, and marks on excluded findings are ignored, following the release rules~\cite{setio2017luna16}. Sensitivity $S(r)$ is measured on each model's own free-response ROC (FROC) at $r$ false positives (FP) per scan. The competition performance metric~\cite{niemeijer2011combining,setio2017luna16} is
\begin{equation}
\mathrm{CPM}=\tfrac{1}{7}\textstyle\sum_{r\in\{\frac18,\frac14,\frac12,1,2,4,8\}} S(r),
\end{equation}
and $S_1=S(1)$. We keep the release's cap of 100 marks per scan and repeat all comparisons without it (supplement Sec.~S3). Because each model reaches each rate on its own FROC, these are \emph{matched} comparisons of ranking quality. A separate analysis applies thresholds fixed on the calibration scans.

\paragraph{Uncertainty.} All effects are paired differences between two models on the same proposals, with 95\% intervals from 2,000 patient-level bootstrap draws~\cite{efron1993bootstrap}. The intervals describe uncertainty over patients for the fitted models, not over training runs. The primary band-versus-ball CPM interval for fit~1 is adjusted for the three comparisons fixed before evaluation (band restoration and the compact head's representation and fusion contrasts); other intervals are pointwise.

\paragraph{Anatomy groups.} Reader contours from LIDC-IDRI~\cite{armato2011lidc,fedorov2020lidc} are rasterized on the native grid, and a voxel is \emph{lesion tissue} if a strict majority of the linked readers include it. Before scores were joined, each lesion with contours and a matching proposal was labelled \emph{reached} if the band of any matching proposal contains a lesion-tissue voxel, and \emph{never reached} otherwise. For each model, the whole proposal stream defines one FP axis; each group's sensitivity is read on this shared axis and divided by the group size. Comparisons at 1~FP/scan therefore use the same global false-positive budget.
